\chapter{Adjoint Functor Theorems}

\begin{overviewbox}
Chapter 7 explained \emph{what} adjunctions are and showed that they are
ubiquitous. But it left open a hard question: given a functor $R$, how do
we know whether it has a left adjoint? In examples we often guess the
left adjoint and verify the bijection. But for many naturally occurring
functors there is no obvious candidate.

This chapter develops the \term{Adjoint Functor Theorems} (AFTs):
existence theorems that say, under suitable conditions, the answer is
yes. The two main statements are:

\begin{itemize}
  \item \term{The General Adjoint Functor Theorem (GAFT)}.\quad
        If $R$ preserves all small limits and satisfies a smallness
        condition (the \emph{solution-set condition}), then $R$ has a
        left adjoint.
  \item \term{The Special Adjoint Functor Theorem (SAFT)}.\quad
        If $R$ preserves all small limits, and its source category is
        \emph{well-powered, complete, and has a small cogenerating set},
        then $R$ has a left adjoint.
\end{itemize}

The dual statements give criteria for right adjoints. Behind both
theorems lies the same idea: a left adjoint $L A$ would have to be the
``initial object'' in the category of $A$-pointed targets; the theorems
provide conditions guaranteeing this initial object exists.

This is the natural climax of the book to date. We have built the cycle
all the way from objects and morphisms through limits, adjunctions, and
Yoneda. The Adjoint Functor Theorems use \emph{all} of this machinery.
\end{overviewbox}


% ═══════════════════════════════════════════════════════════════════════════
\begin{nameorigin}
The \term{Adjoint Functor Theorem} (AFT) was proved by \emph{Peter
Freyd} in his 1964 book \emph{Abelian Categories}. Freyd's theorem
characterized when a functor between locally small categories admits a
left/right adjoint, in terms of two conditions: (1) preserves the
appropriate limits/colimits; (2) satisfies a smallness condition
(the \emph{solution-set condition}). The two main variants are:

\begin{itemize}
  \item \term{GAFT} (General AFT): solution-set condition.
  \item \term{SAFT} (Special AFT): well-powered + cogenerating set.
\end{itemize}

The \term{solution-set condition} is named for the technique it uses:
to find a left adjoint $L A$, one needs a ``solution set'' — a small
collection of arrows out of $A$ that captures every possible target.
The condition trades size for existence, in the standard set-theoretic
manner. Modern formulations (Vol IV Ch.~52, presentable case) replace
solution-set by accessibility, but the spirit is unchanged.
\end{nameorigin}

\section{The Adjoint Functor Question}
\label{sec:aft-question}
% ═══════════════════════════════════════════════════════════════════════════

\subsection{Informally: Existence Without Construction}

When we have an explicit construction of a left adjoint (free monoid,
free group, free abelian group, abelianization, \ldots), we can check
the hom-set bijection directly. But the categorical perspective sometimes
gives us a functor $R$ \emph{first}, and asks whether a left adjoint
\emph{exists}, without telling us how to build one.

For example: ``Does the inclusion of compact Hausdorff spaces into all
topological spaces have a left adjoint?'' Yes --- it is the
Stone--\v{C}ech compactification. But this is a hard theorem to prove
directly; one might want a general criterion that gives existence for
free.

The Adjoint Functor Theorems are exactly such criteria.

\subsection{Expanding: A Reformulation via Initial Objects}

Suppose we want a left adjoint $L A$ for a fixed object $A$ of $\mathcal{C}$.
Define the \term{comma category} $(A \downarrow R)$:
\begin{itemize}
  \item Objects: pairs $(B, f)$ with $B \in \mathcal{D}$ and $f : A \to R B$.
  \item Morphisms $(B, f) \to (B', f')$: morphisms $g : B \to B'$ in
        $\mathcal{D}$ such that $R(g) \circ f = f'$.
\end{itemize}
The object $L A$ together with the unit $\eta_A : A \to R(L A)$ is an
\emph{initial object} of $(A \downarrow R)$: every other pair $(B, f)$
admits a unique morphism from $(L A, \eta_A)$ to itself. So:
\begin{quote}
\textit{Existence of $L A$ $\;\iff\;$ existence of an initial object in
the comma category $(A \downarrow R)$.}
\end{quote}

So the existence of a left adjoint reduces to: do certain comma
categories all have initial objects?

\subsection{Formalizing: The Reduction}

\textbf{Theorem.}\quad A functor $R : \mathcal{D} \to \mathcal{C}$ has a
left adjoint $L$ iff for every $A \in \mathcal{C}$, the comma category
$(A \downarrow R)$ has an initial object. In that case, the initial object
of $(A \downarrow R)$ is $(L A, \eta_A)$, where $\eta$ is the unit.

This is the universal-arrow form of the adjunction (\S 7.1, the third
equivalent definition).

\subsection{Reorganizing: When Does $(A \downarrow R)$ Have an Initial Object?}

In general, having an initial object is a \emph{strong} condition --- many
categories have none. But $(A \downarrow R)$ inherits structure from
$\mathcal{C}$ and $\mathcal{D}$. The strategy of the AFTs is:

\begin{itemize}
  \item \textbf{Get limits in $(A \downarrow R)$.}\quad If $R$ preserves
        limits in $\mathcal{D}$ and $\mathcal{C}$ has the right
        properties, then $(A \downarrow R)$ has small limits.
  \item \textbf{Take a weakly initial set.}\quad A \emph{set} of objects
        from which every other object has a morphism. Limits over the
        weakly initial set yield an initial object.
\end{itemize}

The smallness condition --- existence of a weakly initial set --- is
exactly the \term{solution-set condition} of GAFT.

\subsection{Simplifying: The Strategy in One Sentence}

\textit{To show $R$ has a left adjoint, show $R$ preserves small limits
and that for each $A$ the ``cone of targets'' over $A$ is small.}

\subsection{Formally: Comma Categories Live Above Both Sides}

\formalnote{Comma categories}
The comma category $(A \downarrow R)$ has projection functors to both
$\mathcal{C}$ (constant at $A$, essentially) and $\mathcal{D}$ (sending
$(B, f)$ to $B$). Limits in $(A \downarrow R)$ are computed by taking
limits of the underlying objects and morphisms; the universal property of
the projection makes this work.


% ═══════════════════════════════════════════════════════════════════════════
\section{The General Adjoint Functor Theorem}
\label{sec:gaft}
% ═══════════════════════════════════════════════════════════════════════════

\subsection{Informally: Limits + Smallness $\Rightarrow$ Adjoint}

GAFT is the most general of the adjoint-functor theorems.

\textbf{Theorem (Freyd's General Adjoint Functor Theorem).}\quad Let
$\mathcal{D}$ be a complete category, and $R : \mathcal{D} \to \mathcal{C}$
a functor. Then $R$ has a left adjoint iff:

\begin{enumerate}
  \item $R$ preserves all small limits, and
  \item $R$ satisfies the \term{solution-set condition}: for every
        $A \in \mathcal{C}$, there is a set $I$ of objects
        $\{B_i\}_{i \in I}$ in $\mathcal{D}$ and morphisms
        $f_i : A \to R(B_i)$ such that every morphism $A \to R B$ in
        $\mathcal{C}$ factors as $A \xrightarrow{f_i} R B_i
        \xrightarrow{R(g)} R B$ for some $i$ and some $g : B_i \to B$
        in $\mathcal{D}$.
\end{enumerate}

\subsection{Expanding: Reading the Solution-Set Condition}

The solution-set condition says: for each $A$, there is a \emph{set} (not
a proper class) of ``candidate units'' $f_i : A \to R(B_i)$ such that
every actual unit factors through one of them. This is a
\emph{smallness} condition --- it tames a potentially proper-class-sized
collection of candidates down to a set.

In practice, the solution-set condition holds in any reasonable
category. The condition is automatic if $\mathcal{D}$ has, for instance,
a small dense subcategory or a small cogenerating set.

\subsection{Formalizing: The Statement, Again}

For clarity, restate. Given $R : \mathcal{D} \to \mathcal{C}$ with
$\mathcal{D}$ complete:

\begin{align*}
  & R \text{ has a left adjoint } L \\
  \iff\;
  & R \text{ preserves small limits} \\
  & \quad \text{and for each } A \in \mathcal{C},
    \;\; \exists \text{ set } \{(B_i, f_i)\}_{i \in I} \\
  & \quad \text{such that any } (B, f) \text{ factors through some } (B_i, f_i).
\end{align*}

\textbf{Proof sketch.}\quad ($\Leftarrow$) Build $L A$ as the limit of
the diagram of all factorizations in the comma category through the
solution set. Use limit-preservation of $R$ to verify the universal
property. The solution-set condition is what guarantees this limit is
indexed by a small enough diagram.

($\Rightarrow$) If $L$ exists, then $(L A, \eta_A)$ is the initial object
of $(A \downarrow R)$; take $\{L A\}$ to be the solution set. Limit
preservation is RAPL. \qed

\subsection{Reorganizing: When Each Hypothesis Bites}

\textbf{Completeness of $\mathcal{D}$} is what lets us construct $L A$
as a limit.

\textbf{$R$ preserves limits} is what makes the constructed $L A$ have
the right universal property.

\textbf{Solution-set condition} is what makes the construction
\emph{size-feasible}: without it, the relevant ``limit'' would be over a
proper class.

\subsection{Simplifying: GAFT as a Recipe}

To verify GAFT:
\begin{enumerate}
  \item Check $\mathcal{D}$ has all small limits.
  \item Check $R$ preserves them.
  \item Check the solution-set condition (often the trickiest step).
\end{enumerate}

\subsection{Formally: The Dual for Right Adjoints}

\formalnote{GAFT for right adjoints}
By duality: a functor $L : \mathcal{C} \to \mathcal{D}$ with
$\mathcal{D}$ cocomplete has a right adjoint iff $L$ preserves all small
colimits and satisfies the dual solution-set condition.


% ═══════════════════════════════════════════════════════════════════════════
\section{The Special Adjoint Functor Theorem}
\label{sec:saft}
% ═══════════════════════════════════════════════════════════════════════════

\subsection{Informally: A Cleaner Statement Under Extra Conditions}

GAFT is general but requires checking the solution-set condition, which
is sometimes awkward. SAFT replaces the solution-set condition with two
cleaner structural assumptions on $\mathcal{D}$: \emph{well-poweredness}
and \emph{a small cogenerating set}.

\subsection{Expanding: Well-Powered and Cogenerating}

\textbf{Well-powered.}\quad A category $\mathcal{D}$ is
\term{well-powered} if for every object $B$, the collection of
subobjects of $B$ (equivalence classes of monomorphisms into $B$) is a
set, not a proper class. This holds in almost every category one meets
in practice.

\textbf{Cogenerating set.}\quad A set $\{S_j\}_{j \in J}$ in $\mathcal{D}$
is \term{cogenerating} (or a ``set of cogenerators'') if the functor
\begin{equation*}
  \mathcal{D} \to \prod_j \mathbf{Set},
  \quad
  B \mapsto (\mathrm{Hom}(B, S_j))_{j \in J}
\end{equation*}
is faithful; equivalently, morphisms $f, g : B \to B'$ are equal iff
$h \circ f = h \circ g$ for every $h : B' \to S_j$ for every $j$.

\subsection{Formalizing: SAFT}

\textbf{Theorem (Special Adjoint Functor Theorem).}\quad Let
$\mathcal{D}$ be complete, well-powered, and have a small cogenerating
set. Let $R : \mathcal{D} \to \mathcal{C}$ be a functor. Then $R$ has a
left adjoint iff $R$ preserves all small limits.

\subsection{Reorganizing: Why SAFT Is Easier}

The two structural conditions on $\mathcal{D}$ (well-powered and
cogenerating set) together imply the solution-set condition of GAFT for
every $R$. So under SAFT's hypotheses, you do not need to verify
solution-sets separately; you only check limit-preservation.

In categories like $\mathbf{CHaus}$ (compact Hausdorff spaces) or any
locally presentable category, SAFT's hypotheses are automatic, so the
existence of an adjoint reduces to checking limit-preservation.

\subsection{Simplifying: SAFT in Practice}

For ``well-behaved'' target categories: limit-preservation is enough.

\subsection{Formally: Locally Presentable Categories}

\formalnote{A modern reformulation}
The most powerful adjoint-functor theorem in current use is the AFT for
\term{locally presentable categories}: in a locally presentable category,
a functor between locally presentable categories has a left adjoint iff
it preserves limits. This subsumes both GAFT and SAFT for an enormous
class of practical examples.


% ═══════════════════════════════════════════════════════════════════════════
\section{Worked Examples}
\label{sec:aft-examples}
% ═══════════════════════════════════════════════════════════════════════════

\subsection{Informally: Watching the Theorems Work}

We apply the theorems to three classical adjoints that we have seen
existing for ``other'' reasons, and observe the AFT machinery at work.

\subsection{Expanding: Three Workouts}

\begin{example}{Stone--\v{C}ech Compactification, via SAFT}
The inclusion $i : \mathbf{CHaus} \hookrightarrow \mathbf{Top}$ has a
left adjoint $\beta$ (the Stone--\v{C}ech compactification). To prove
existence by SAFT, we need:
\begin{itemize}
  \item $\mathbf{CHaus}$ is complete (yes; limits of compact Hausdorff
        spaces are compact Hausdorff).
  \item $\mathbf{CHaus}$ is well-powered (yes; subobjects of a compact
        Hausdorff space are closed subsets, which form a set).
  \item $\mathbf{CHaus}$ has a small cogenerating set: $\{[0, 1]\}$ is
        one. (Two morphisms into a $\mathrm{CHaus}$ space agree iff
        they agree after composition with every map to $[0, 1]$, which
        is Urysohn-like.)
  \item $i$ preserves limits (yes; it preserves the underlying topology).
\end{itemize}
SAFT delivers the adjoint $\beta$ without our needing to construct it
explicitly.
\end{example}

\begin{example}{Free $R$-Module Functor, via GAFT}
Let $R$ be a ring. The forgetful $U : R\text{-}\mathbf{Mod} \to \mathbf{Set}$
has a left adjoint (the free $R$-module functor). $R\text{-}\mathbf{Mod}$
is complete, $U$ preserves limits, and the solution-set condition holds
with a singleton set: $\{R^X\}$, the free module on the set $X$. GAFT
delivers existence.
\end{example}

\begin{example}{Sheafification, via SAFT in $[\mathcal{C}^{\mathrm{op}}, \mathbf{Set}]$}
The inclusion of sheaves into presheaves is a fully faithful functor with
a left adjoint, called \term{sheafification}. SAFT applies: the presheaf
category is complete, well-powered, and the family of representables
$\{Y(C)\}_{C \in \mathcal{C}}$ is jointly cogenerating. The inclusion of
sheaves preserves limits because limits of sheaves are computed
pointwise. SAFT gives sheafification as a left adjoint.
\end{example}

\subsection{Formalizing: A General Recipe}

When confronted with ``does $R$ have a left adjoint?'', try SAFT first
(it requires less work to verify). If $\mathcal{D}$ is not known to be
well-powered or cogenerated, fall back to GAFT and check the
solution-set condition by hand.

If even GAFT fails, there is a fundamental obstruction: either $R$ does
not preserve limits (and so no left adjoint exists at all), or the
candidate left adjoint involves a proper-class-sized colimit and is not
representable.

\subsection{Reorganizing: When Adjoints Fail to Exist}

\begin{itemize}
  \item The inclusion $\mathbf{Set} \hookrightarrow \mathbf{Class}$ has
        no left adjoint, because $\mathbf{Class}$ is not a category in
        the usual sense (size issues).
  \item The forgetful $\mathbf{TopGrp} \to \mathbf{Grp}$ (forgetting the
        topology of a topological group) has no left adjoint, because
        the candidate would be a free topological group on a discrete
        group, which exists but is hard to verify without further
        smallness.
\end{itemize}

\subsection{Simplifying: The Final Slogan}

\begin{equation*}
  \boxed{\;\;
    \text{Preserves all limits} + \text{Smallness}
    \;\Rightarrow\;
    \text{Has a left adjoint}.
  \;\;}
\end{equation*}

\subsection{Formally: Connection to Yoneda}

\formalnote{AFT and representability}
The Adjoint Functor Theorems are really representability theorems in
disguise. Asking whether $R$ has a left adjoint is asking whether each
of the presheaves $\mathrm{Hom}_\mathcal{C}(A, R(-))$ is representable.
The theorems give conditions under which a presheaf is representable, and
left-adjoint-existence is a special case.


% ═══════════════════════════════════════════════════════════════════════════
\section*{Exercises}
\addcontentsline{toc}{section}{Exercises}
% ═══════════════════════════════════════════════════════════════════════════

\begin{exercisesbox}
\begin{qa}
\qaitem{What question does an Adjoint Functor Theorem answer?}{Given a functor $R$, does it have a left adjoint?}
\qaitem{When does $R$ certainly \emph{not} have a left adjoint?}{When $R$ fails to preserve some small limit.}
\qaitem{Why?}{By RAPL: right adjoints must preserve all small limits.}
\qaitem{What is the comma category $(A \downarrow R)$?}{Pairs $(B, f : A \to R B)$, with the obvious morphisms.}
\qaitem{Why does $(A \downarrow R)$ matter for adjunctions?}{$L A$ exists iff $(A \downarrow R)$ has an initial object.}
\qaitem{What initial object does it have?}{$(L A, \eta_A)$, the universal arrow from $A$ to $R$.}
\qaitem{State GAFT in slogan form.}{Complete target, $R$ preserves limits, plus solution-set $\Rightarrow$ $R$ has a left adjoint.}
\qaitem{What is the solution-set condition?}{For each $A$, a set of $f_i : A \to R B_i$ such that every $f : A \to R B$ factors through one.}
\qaitem{Why does the condition need a \emph{set}?}{To avoid proper-class-sized constructions.}
\qaitem{In the SAFT, which two extra assumptions replace the solution-set condition?}{Well-poweredness and a small cogenerating set.}
\qaitem{What is a well-powered category?}{One in which each object has only a set's worth of subobjects.}
\qaitem{What is a cogenerating set?}{A set $\{S_j\}$ such that hom-into-cogenerators is jointly faithful.}
\qaitem{Give an example of a cogenerator in $\mathbf{CHaus}$.}{$[0, 1]$.}
\qaitem{Where does SAFT apply that GAFT doesn't need to be invoked?}{When the target category is locally presentable; SAFT often suffices.}
\qaitem{What is the connection to Yoneda?}{Both GAFT and SAFT are representability theorems for the presheaves $\mathrm{Hom}(A, R(-))$.}
\qaitem{What strategy underlies both theorems?}{Construct $L A$ as a limit over the comma category, controlled by a smallness hypothesis.}
\qaitem{Why does completeness of $\mathcal{D}$ enter?}{Because we construct $L A$ as a limit of objects in $\mathcal{D}$.}
\qaitem{Why does limit-preservation of $R$ enter?}{Because $L A$'s universal property must hold after applying $R$.}
\qaitem{Where does the solution-set condition enter the proof?}{To ensure the limit we take is over a set-sized diagram.}
\qaitem{Can a functor have neither a left nor a right adjoint?}{Yes; many do.}
\qaitem{Can a functor have both?}{Yes; e.g., the identity functor has itself on both sides.}
\qaitem{What does ``$f^*$ has both adjoints'' say in logic?}{There are both existential and universal quantifiers, given by left and right adjoints.}
\qaitem{Is the forgetful functor $\mathbf{Mon} \to \mathbf{Set}$ adjoint-bearing?}{Yes; it has a left adjoint (free monoid), no right adjoint.}
\qaitem{Why no right adjoint to forgetful $\mathbf{Mon} \to \mathbf{Set}$?}{The forgetful doesn't preserve coproducts on the nose.}
\qaitem{In a locally presentable category, what's the criterion for an adjoint?}{Limit preservation suffices.}
\qaitem{What's the dual of GAFT for right adjoints?}{Cocomplete target, $L$ preserves colimits, plus dual solution-set condition.}
\qaitem{Sheafification is an example of which theorem applied?}{SAFT.}
\qaitem{Why is the AFT considered a deep theorem?}{It turns existence-of-adjoint into a checkable list of conditions, often the only practical route.}
\qaitem{What's the central tool common to all the theorems?}{Comma categories with initial objects.}
\qaitem{Final slogan?}{``Preserves all limits + smallness $\Rightarrow$ has a left adjoint.''}
\end{qa}
\end{exercisesbox}


% ═══════════════════════════════════════════════════════════════════════════
\section*{Chapter Summary}
\addcontentsline{toc}{section}{Chapter Summary}
% ═══════════════════════════════════════════════════════════════════════════

\begin{summarybox}
The \textbf{Adjoint Functor Theorems} give criteria for a functor $R$ to
have a left adjoint.

\textbf{GAFT}: If $\mathcal{D}$ is complete, $R$ preserves small limits,
and $R$ satisfies the \textbf{solution-set condition}, then $R$ has a
left adjoint.

\textbf{SAFT}: If additionally $\mathcal{D}$ is well-powered and has a
small cogenerating set, the solution-set condition is automatic; limit
preservation suffices.

The underlying mechanism: $L A$ exists iff the comma category
$(A \downarrow R)$ has an initial object, which it does under the
hypotheses, constructed as a limit over a solution-set-indexed diagram.

These theorems supply existence of Stone--\v{C}ech compactification,
sheafification, free module functors, and many other classical
left adjoints --- without requiring an explicit construction in each case.
\end{summarybox}


% ═══════════════════════════════════════════════════════════════════════════
\section*{Formal Restatement}
\addcontentsline{toc}{section}{Formal Restatement}
% ═══════════════════════════════════════════════════════════════════════════

\begin{formalrestatement}
\textbf{Comma category.}\quad
$(A \downarrow R)$ has objects $(B, f)$ with $f : A \to R B$, and
morphisms $g : B \to B'$ with $R(g) \circ f = f'$.

\medskip
\textbf{Adjoint criterion.}\quad
$R$ has a left adjoint $L$ iff $(A \downarrow R)$ has an initial object
for every $A$. That initial object is $(L A, \eta_A)$.

\medskip
\textbf{GAFT.}\quad
If $\mathcal{D}$ is complete, then $R : \mathcal{D} \to \mathcal{C}$ has
a left adjoint iff:
\begin{itemize}
  \item $R$ preserves all small limits, and
  \item for every $A$, there is a set $\{(B_i, f_i)\}$ such that every
        $f : A \to R B$ factors as $R(g) \circ f_i$ for some $i$,
        $g : B_i \to B$.
\end{itemize}

\medskip
\textbf{SAFT.}\quad
If $\mathcal{D}$ is complete, well-powered, and has a small cogenerating
set, then $R : \mathcal{D} \to \mathcal{C}$ has a left adjoint iff $R$
preserves all small limits.

\medskip
\textbf{Dual.}\quad
Apply to $\mathcal{D}^{\mathrm{op}}$ for right-adjoint existence.

\medskip
\noindent\textit{This concludes the foundational arc of the book.
Chapter~11 onward will turn to connections with other areas of
mathematics, where the toolkit of categories, functors, natural
transformations, limits, adjunctions, monads, Yoneda, and AFT becomes
the working vocabulary.}
\end{formalrestatement}
